\section{Workload and electricity-price inputs}
\subsection{Adopted workload arrays and their interpretation}
Published studies describe differing utilisation levels, workload classes and temporal flexibility across facilities~\cite{Cao2022ApEn,6779441,Misaghian2023TIA,Liu2012PER}. They motivate separating immediate work from work that may be rescheduled, but do not identify a universal flexible/inflexible ratio. The inputs used here are stylised profiles selected with reference to that literature. They are not claimed to be a measured site trace, a statistically representative sample of facilities or a deterministic reconstruction of a particular published trace.

The authoritative numerical inputs are \path{data/load_profiles.csv} and \path{data/shiftability_profile.csv}, whose active slots 1--96 are reproduced jointly in Table~\ref{tab:workload-inputs}. The load file also contains legacy rows 97--108; annual timeline construction does not use those rows. The former file specifies flexible and inflexible arrival utilisation; the latter specifies four fractions summing to one in each active slot. These exact quarter-hour values supersede abbreviated hourly summaries when reproducing the calculations. In particular, the tranche array is used at its recorded quarter-hour resolution; no hourly table is expanded to replace it. The main paper retains its earlier literature-range table as context, with the adopted columns updated to hourly means. Those range summaries do not replace the exact quarter-hour inputs.

For an arrival in interval $t$, class $k$ receives work
\begin{equation}
a_{t,k}=u_t^{\mathrm{flex,arr}}\alpha_{t,k}\Delta t,
\quad \ell_{t,k}=t+D_k\Delta t,
\quad (D_1,D_2,D_3,D_4)=(2,4,8,12).
\end{equation}
The final execution interval starts at $\ell_{t,k}$ and closes at $\ell_{t,k}+\Delta t$. For example, a three-hour class arriving at 12:00 may execute in the interval starting at 15:00 and ending at 15:15. Thus the class names denote shifts between interval starts, not a guarantee of completion within exactly three continuous hours of an interval's opening. Cohort creation omits quantities no larger than the $10^{-7}$ CPU-hour numerical tolerance.

For workload multiplier $m$, let $v_t=u_t^{\mathrm{inflex}}+u_t^{\mathrm{flex,arr}}$ in the original input. The sensitivity uses
\begin{equation}
u_t^{\mathrm{flex,arr}}(m)=\min\{m u_t^{\mathrm{flex,arr}}(1),v_t\},
\qquad u_t^{\mathrm{inflex}}(m)=v_t-u_t^{\mathrm{flex,arr}}(m).
\end{equation}
Tranche fractions remain fixed. The clipping rule preserves total arriving CPU-hours and avoids creating more flexible work than the total workload. A multiplier is therefore not always an exactly proportional increase in the realised flexible share.

\begin{longtable}{@{}lrrrrrr@{}}
\caption{Actual quarter-hour workload inputs. Utilisations are fractions of rated CPU capacity; tranche shares are fractions of flexible work.}\label{tab:workload-inputs}\\\toprule
Local start & Inflexible & Flexible & $\alpha_1$ & $\alpha_2$ & $\alpha_3$ & $\alpha_4$ \\\midrule\endfirsthead
\toprule
Local start & Inflexible & Flexible & $\alpha_1$ & $\alpha_2$ & $\alpha_3$ & $\alpha_4$ \\\midrule\endhead
00:00 & 0.28 & 0.40 & 0.25 & 0.25 & 0.20 & 0.30 \\
00:15 & 0.27 & 0.38 & 0.25 & 0.25 & 0.20 & 0.30 \\
00:30 & 0.26 & 0.36 & 0.25 & 0.25 & 0.20 & 0.30 \\
00:45 & 0.25 & 0.34 & 0.25 & 0.25 & 0.20 & 0.30 \\
01:00 & 0.25 & 0.31 & 0.25 & 0.25 & 0.15 & 0.35 \\
01:15 & 0.23 & 0.32 & 0.25 & 0.25 & 0.15 & 0.35 \\
01:30 & 0.21 & 0.32 & 0.25 & 0.25 & 0.15 & 0.35 \\
01:45 & 0.19 & 0.32 & 0.25 & 0.25 & 0.15 & 0.35 \\
02:00 & 0.17 & 0.33 & 0.35 & 0.17 & 0.18 & 0.30 \\
02:15 & 0.16 & 0.30 & 0.35 & 0.17 & 0.18 & 0.30 \\
02:30 & 0.16 & 0.28 & 0.35 & 0.17 & 0.18 & 0.30 \\
02:45 & 0.16 & 0.26 & 0.35 & 0.17 & 0.18 & 0.30 \\
03:00 & 0.16 & 0.23 & 0.25 & 0.15 & 0.20 & 0.40 \\
03:15 & 0.14 & 0.24 & 0.25 & 0.15 & 0.20 & 0.40 \\
03:30 & 0.12 & 0.25 & 0.25 & 0.15 & 0.20 & 0.40 \\
03:45 & 0.10 & 0.26 & 0.25 & 0.15 & 0.20 & 0.40 \\
04:00 & 0.08 & 0.27 & 0.25 & 0.20 & 0.15 & 0.40 \\
04:15 & 0.07 & 0.27 & 0.25 & 0.20 & 0.15 & 0.40 \\
04:30 & 0.06 & 0.27 & 0.25 & 0.20 & 0.15 & 0.40 \\
04:45 & 0.06 & 0.27 & 0.25 & 0.20 & 0.15 & 0.40 \\
05:00 & 0.06 & 0.27 & 0.25 & 0.15 & 0.15 & 0.45 \\
05:15 & 0.07 & 0.24 & 0.25 & 0.15 & 0.15 & 0.45 \\
05:30 & 0.09 & 0.22 & 0.25 & 0.15 & 0.15 & 0.45 \\
05:45 & 0.10 & 0.20 & 0.25 & 0.15 & 0.15 & 0.45 \\
06:00 & 0.12 & 0.18 & 0.23 & 0.28 & 0.15 & 0.34 \\
06:15 & 0.14 & 0.19 & 0.23 & 0.28 & 0.15 & 0.34 \\
06:30 & 0.16 & 0.21 & 0.23 & 0.28 & 0.15 & 0.34 \\
06:45 & 0.18 & 0.22 & 0.23 & 0.28 & 0.15 & 0.34 \\
07:00 & 0.20 & 0.24 & 0.20 & 0.20 & 0.15 & 0.45 \\
07:15 & 0.21 & 0.24 & 0.20 & 0.20 & 0.15 & 0.45 \\
07:30 & 0.22 & 0.24 & 0.20 & 0.20 & 0.15 & 0.45 \\
07:45 & 0.23 & 0.24 & 0.20 & 0.20 & 0.15 & 0.45 \\
08:00 & 0.24 & 0.24 & 0.25 & 0.15 & 0.13 & 0.47 \\
08:15 & 0.26 & 0.25 & 0.25 & 0.15 & 0.13 & 0.47 \\
08:30 & 0.29 & 0.26 & 0.25 & 0.15 & 0.13 & 0.47 \\
08:45 & 0.31 & 0.27 & 0.25 & 0.15 & 0.13 & 0.47 \\
09:00 & 0.34 & 0.28 & 0.25 & 0.17 & 0.16 & 0.42 \\
09:15 & 0.35 & 0.25 & 0.25 & 0.17 & 0.16 & 0.42 \\
09:30 & 0.36 & 0.23 & 0.25 & 0.17 & 0.16 & 0.42 \\
09:45 & 0.37 & 0.21 & 0.25 & 0.17 & 0.16 & 0.42 \\
10:00 & 0.37 & 0.19 & 0.35 & 0.22 & 0.13 & 0.30 \\
10:15 & 0.38 & 0.18 & 0.35 & 0.22 & 0.13 & 0.30 \\
10:30 & 0.39 & 0.17 & 0.35 & 0.22 & 0.13 & 0.30 \\
10:45 & 0.40 & 0.16 & 0.35 & 0.22 & 0.13 & 0.30 \\
11:00 & 0.42 & 0.16 & 0.32 & 0.20 & 0.15 & 0.33 \\
11:15 & 0.41 & 0.17 & 0.32 & 0.20 & 0.15 & 0.33 \\
11:30 & 0.40 & 0.18 & 0.32 & 0.20 & 0.15 & 0.33 \\
11:45 & 0.39 & 0.19 & 0.32 & 0.20 & 0.15 & 0.33 \\
12:00 & 0.40 & 0.20 & 0.45 & 0.25 & 0.15 & 0.15 \\
12:15 & 0.39 & 0.21 & 0.45 & 0.25 & 0.15 & 0.15 \\
12:30 & 0.39 & 0.22 & 0.45 & 0.25 & 0.15 & 0.15 \\
12:45 & 0.38 & 0.23 & 0.45 & 0.25 & 0.15 & 0.15 \\
13:00 & 0.36 & 0.24 & 0.50 & 0.20 & 0.15 & 0.15 \\
13:15 & 0.36 & 0.24 & 0.50 & 0.20 & 0.15 & 0.15 \\
13:30 & 0.35 & 0.25 & 0.50 & 0.20 & 0.15 & 0.15 \\
13:45 & 0.35 & 0.26 & 0.50 & 0.20 & 0.15 & 0.15 \\
14:00 & 0.35 & 0.27 & 0.40 & 0.20 & 0.25 & 0.15 \\
14:15 & 0.34 & 0.27 & 0.40 & 0.20 & 0.25 & 0.15 \\
14:30 & 0.33 & 0.27 & 0.40 & 0.20 & 0.25 & 0.15 \\
14:45 & 0.32 & 0.27 & 0.40 & 0.20 & 0.25 & 0.15 \\
15:00 & 0.33 & 0.27 & 0.45 & 0.25 & 0.15 & 0.15 \\
15:15 & 0.34 & 0.25 & 0.45 & 0.25 & 0.15 & 0.15 \\
15:30 & 0.35 & 0.23 & 0.45 & 0.25 & 0.15 & 0.15 \\
15:45 & 0.37 & 0.21 & 0.45 & 0.25 & 0.15 & 0.15 \\
16:00 & 0.40 & 0.20 & 0.50 & 0.15 & 0.20 & 0.15 \\
16:15 & 0.37 & 0.25 & 0.50 & 0.15 & 0.20 & 0.15 \\
16:30 & 0.33 & 0.30 & 0.50 & 0.15 & 0.20 & 0.15 \\
16:45 & 0.30 & 0.35 & 0.50 & 0.15 & 0.20 & 0.15 \\
17:00 & 0.27 & 0.40 & 0.50 & 0.20 & 0.15 & 0.15 \\
17:15 & 0.27 & 0.41 & 0.50 & 0.20 & 0.15 & 0.15 \\
17:30 & 0.27 & 0.42 & 0.50 & 0.20 & 0.15 & 0.15 \\
17:45 & 0.26 & 0.43 & 0.50 & 0.20 & 0.15 & 0.15 \\
18:00 & 0.26 & 0.45 & 0.10 & 0.20 & 0.30 & 0.40 \\
18:15 & 0.26 & 0.45 & 0.10 & 0.20 & 0.30 & 0.40 \\
18:30 & 0.26 & 0.45 & 0.10 & 0.20 & 0.30 & 0.40 \\
18:45 & 0.26 & 0.45 & 0.10 & 0.20 & 0.30 & 0.40 \\
19:00 & 0.26 & 0.45 & 0.15 & 0.20 & 0.25 & 0.40 \\
19:15 & 0.26 & 0.45 & 0.15 & 0.20 & 0.25 & 0.40 \\
19:30 & 0.25 & 0.46 & 0.15 & 0.20 & 0.25 & 0.40 \\
19:45 & 0.25 & 0.46 & 0.15 & 0.20 & 0.25 & 0.40 \\
20:00 & 0.25 & 0.47 & 0.18 & 0.12 & 0.25 & 0.45 \\
20:15 & 0.26 & 0.46 & 0.18 & 0.12 & 0.25 & 0.45 \\
20:30 & 0.27 & 0.45 & 0.18 & 0.12 & 0.25 & 0.45 \\
20:45 & 0.28 & 0.44 & 0.18 & 0.12 & 0.25 & 0.45 \\
21:00 & 0.29 & 0.41 & 0.22 & 0.18 & 0.25 & 0.35 \\
21:15 & 0.29 & 0.41 & 0.22 & 0.18 & 0.25 & 0.35 \\
21:30 & 0.29 & 0.41 & 0.22 & 0.18 & 0.25 & 0.35 \\
21:45 & 0.30 & 0.41 & 0.22 & 0.18 & 0.25 & 0.35 \\
22:00 & 0.30 & 0.40 & 0.20 & 0.16 & 0.24 & 0.40 \\
22:15 & 0.27 & 0.40 & 0.20 & 0.16 & 0.24 & 0.40 \\
22:30 & 0.25 & 0.41 & 0.20 & 0.16 & 0.24 & 0.40 \\
22:45 & 0.22 & 0.41 & 0.20 & 0.16 & 0.24 & 0.40 \\
23:00 & 0.21 & 0.42 & 0.15 & 0.20 & 0.20 & 0.45 \\
23:15 & 0.21 & 0.42 & 0.15 & 0.20 & 0.20 & 0.45 \\
23:30 & 0.21 & 0.42 & 0.15 & 0.20 & 0.20 & 0.45 \\
23:45 & 0.21 & 0.42 & 0.15 & 0.20 & 0.20 & 0.45 \\
\bottomrule\end{longtable}


\subsection{Prices, calendar and settlement}
The source is LCCC's IMRP actuals dataset~\cite{lccc_imrp_actuals}, whose reference product is the Great Britain day-ahead hourly price. Source fields identify local date, chronological hourly settlement period and IMRP amount in GBP/MWh. The 2025 series covers 8,760 physical hours. Hourly prices are replicated to four consecutive quarter-hours without smoothing, giving 35,040 committed intervals. The timeline uses UTC for physical ordering and Europe/London for local-date and workload-slot labels. Spring and autumn clock changes produce 92 and 100 quarter-hours in those local days; workload values for skipped/repeated local-clock periods are correspondingly skipped/repeated. No physical UTC interval is missing or duplicated.

Negative prices remain signed in both optimisation and reported costs. There are 672 negative-price quarter-hours in the retained annual data. IMRP is a reference-price signal, not a full retail tariff. The first three hourly prices on 1 January 2026 are GBP~46.49, 40.19 and 28.99/MWh; they supply the twelve final look-ahead intervals. The exact expanded central price timeline is included as \path{data/price_timeline.csv}. Annual cost uses only committed intervals in calendar 2025.
